\scanpage{236}
- \mspage{188} - \margin{IX}
\[
\underline R_U/\underline R_U^{\mathrm f}\simeq M_U\ ,\tag{13.5.12}
\]
et l'homomorphisme $i$ de (13.3.4) peut s'écrire comme un homomorphisme
\[
i:M_U\longrightarrow A^{\unsure{\natural|\sharp}}\ ,\tag{13.5.13}
\]
qu'on retrouvera dans un contexte un peu différent au par. suivant.

14. \uline{Liens avec la théorie transcendante} : \uline{cas rigide-analytique}. \uline{L'homomorphisme canonique} $\underline M_K\longrightarrow A_K^{\unsure{\natural|\sharp}0}$.

14.1. Nous supposons connu ici les fondements de la théorie des espaces rigide analytiques, pour lesquels on pourra consulter [12] [13] [34].

Soit $S$ un trait complet, $A_K$ un schéma abélien sur son corps des fractions. On a construit alors dans 8.1, par descente à partir de l'extension de Raynaud dans le cas semi-stable, une extension $G_K=A_K^{\unsure{\natural|\sharp}0}$ (8.1.1)
\[
\text{(14.1.1)}\qquad
\diagram{\begin{tikzcd}0\arrow[r]&T_K\arrow[r]&G_K\arrow[r]&B_K\arrow[r]&0\end{tikzcd}}
\]
d'un schéma abélien $B_K$ par un tore $T_K$. Considérant le schéma abélien dual $A'_K$, on a de même une extension $G'_K=A_K'{}^{\unsure{\natural|\sharp}0}$
\[
\text{(14.1.2)}\qquad
\diagram{\begin{tikzcd}0\arrow[r]&T'_K\arrow[r]&G'_K\arrow[r]&B'_K\arrow[r]&0\end{tikzcd}}\ ,
\]
où $B'_K$ s'identifie d'ailleurs au schéma abélien dual de $B_K$ (8.2). On désigne, comme dans (8.1.8), par $\underline M_K$ et $\underline M'_K$ les réseaux sur $K$ duals de $T'_K$ et de $T_K$ respectivement :
\[
\underline M_K=D(T'_K)\ ,\qquad \underline M'_K=D(T_K)\ .\tag{14.1.3}
\]
\margin{235}
\margin{506}
